\documentclass[letterpaper,11pt]{article}
\usepackage{graphicx}
\usepackage{geometry}
\linespread{1.2}                        
\setlength{\oddsidemargin}{0.0in}            
\setlength{\textwidth}{6.5in}    
\usepackage[utf8]{inputenc} % Asegúrate de usar utf8
\usepackage[T1]{fontenc} % Asegúrate de usar T1 font encoding      
\setlength{\topmargin}{-.6in}            
\setlength{\textheight}{9in}
\usepackage{latexsym} 
\usepackage{amssymb} 
\usepackage{amsmath}
\usepackage{amsxtra}
\usepackage[mathcal]{eucal} 
\usepackage{enumerate}
\usepackage{hyperref}
\usepackage{rotating}
\usepackage{caption}
\usepackage{lscape}
\usepackage{paralist} 
\usepackage{indentfirst}
\usepackage[dvipsnames,svgnames]{xcolor} 
\usepackage{setspace}
\usepackage{multicol} 
\usepackage{float}

\definecolor{codegreen}{rgb}{0,0.6,0}
\definecolor{codegray}{rgb}{0.5,0.5,0.5}
\definecolor{codepurple}{cmyk}{0.65, 1, 0, 0.35}
\definecolor{backcolour}{rgb}{0.95,0.95,0.92}

\usepackage{listings}
\usepackage{xcolor}

\lstset{
    language=R, % lenguaje de programación
    basicstyle=\small\ttfamily, % estilo básico
    commentstyle=\color{gray}, % estilo para los comentarios
    keywordstyle=\color{blue}, % estilo para las palabras clave
    stringstyle=\color{black}, % estilo para las cadenas de texto
    showstringspaces=false, % no mostrar espacios en cadenas de texto
    tabsize=4, % tamaño de la tabulación
    breaklines=true, % permitir el salto de línea automático
    postbreak=\mbox{\textcolor{red}{$\hookrightarrow$}\space}, % símbolo al final de una línea partida
    numbers=left, % mostrar números de línea
    numberstyle=\tiny\color{gray}, % estilo para los números de línea
    frame=single, % mostrar un borde alrededor del código
    backgroundcolor=\color{gray!10}, % color de fondo
}

\begin{document}

\begin{center}
    {\Large  Taller 3: Econometría }

Profesor:    Jocelyn Tapia\\
Ayudante: Jociney Honorio y Matías Sandoval 
\end{center}

\textbf{Instrucciones:} 

\begin{itemize}
    \item Usted dispone de hasta el día 16 de junio a las 22:00  para entregar su taller.
    \item Puede realizarlo de manera individual o en parejas.
       \item Puede utilizar su computador y los apuntes tanto de clase como de ayudantía.
    \item Debe enviar un archivo en formato PDF y un R script (extensión .R)
\end{itemize}



Para responder las preguntas siguientes, utilice la base de datos \textit{htv} 
disponible en la librería \textbf{wooldridge}. Cargue la base eliminando 
observaciones con valores nulos antes de comenzar el análisis.

\begin{enumerate}

\item (10 puntos) Construya una matriz de correlación entre las variables de 
la base de datos. Identifique y comente los pares de variables con correlaciones 
altas (en valor absoluto superiores a 0,7), señalando qué problemas podrían 
generar en la estimación de un modelo de regresión lineal múltiple.

\item (15 puntos) Estime un modelo de regresión lineal múltiple que explique 
el salario de los individuos. Justifique la elección de cada regresor incluido, 
basándose en criterios económicos o en los resultados de la pregunta anterior.

\item (15 puntos) Calcule el Factor de Inflación de Varianza (VIF) para los 
regresores del modelo estimado en la pregunta anterior. Compare los resultados 
con los obtenidos en la matriz de correlación y discuta si existe evidencia de 
multicolinealidad. En caso de detectarla, indique qué variables deberían 
eliminarse del modelo y por qué.

\item (40 puntos) Reestime el modelo incorporando los ajustes definidos en la 
pregunta anterior. A continuación:
\begin{enumerate}[a)]
    \item Aplique el test de Breusch-Pagan y el test de White para evaluar la 
    presencia de heterocedasticidad. Plantee claramente las hipótesis nula y 
    alternativa, y concluya.
    \item Si detecta heterocedasticidad, reestime el modelo utilizando errores 
    estándar robustos (White). Compare los errores estándar obtenidos con y sin 
    corrección y comente las diferencias.
\end{enumerate}

\item (20 puntos) A partir del modelo final estimado, analice la significancia 
estadística de cada coeficiente al 5\% de significancia. Identifique qué 
variables tienen un efecto estadísticamente significativo sobre el salario e 
interprete el signo y magnitud de sus coeficientes.

\end{enumerate}


\end{document}
